% GENERATED FILE. Edit canonical lessons, then run npm run generate:books.
% canonical-source-hash: fnv1a64:de27cb64d5c7caf3
% canonical-lessons: HI-C390-laksya, HI-C390-iccha, HI-C390-sankalp, HI-C390-qadam, HI-C390-kosiskarna, HI-R390-review-383-390

\chapter{Setting a Goal}
\label{ch:hi-setting-a-goal}
\hlchaptermodality{390}

\begin{chapteropening}
\textbf{By the end of this chapter:} \emph{I can say what I want to reach and what it will take.}

Complete the last lesson of chapter 390: I can say what I want to reach and what it will take.
\end{chapteropening}

\section[\texorpdfstring{\hi{लक्ष्य} (lakṣya)}{लक्ष्य (lakṣya)}]{\hi{लक्ष्य} (lakṣya) — a goal, a target}

\label{lesson:HI-C390-laksya}



\begin{quote}
Before the new one: say the Hindi for normal, then the Hindi for to brighten.
\end{quote}

\subsection*{You'll want to know: \hi{लक्ष्य}}
\textbf{\hi{लक्ष्य}} — \emph{lakṣya} — \textquotedblleft{}a goal, a target\textquotedblright{}.

\begin{quote}
\hi{अगर} \hi{तुम्हारा} \hi{लक्ष्य} \hi{साफ़} \hi{होगा} \hi{तो} \hi{रास्ता} \hi{आसान} \hi{होगा।}
\end{quote}

\begin{quote}
\emph{agar tumhārā lakṣya sāf hogā to rāstā āsān hogā.}
\end{quote}

\begin{quote}
If your goal is clear, the way will be easy.
\end{quote}

Masculine.

\begin{cousinweb}
\textbf{\hi{लक्ष्य}} (\emph{lakṣya}) is a tatsama, borrowed from Sanskrit \emph{lakṣya}, \textquotedblleft{}what is kept in view; a mark\textquotedblright{} (Monier-Williams).
\end{cousinweb}

\begin{grammarlens}[title={What you've built}]
One more word for saying what happens if.
\end{grammarlens}

\subsection*{Your turn}
\begin{itemize}
\raggedright
  \item \emph{Say it:} \emph{lakṣya}
  \item \emph{Say it:} \emph{lakṣya}, once more
  \item \emph{Say it:} \emph{agar tumhārā lakṣya sāf hogā to rāstā āsān hogā.} — If your goal is clear, the way will be easy.
  \item \emph{Recall:} say the Hindi for normal, then the Hindi for to brighten, then say \emph{lakṣya} again
\end{itemize}

\begin{tcolorbox}[breakable,colback=teal!4,colframe=teal!35!black,title={Before you move on}]
What does \hi{लक्ष्य} mean? (\textquotedblleft{}A goal, a target\textquotedblright{}.) Say it once more.
\end{tcolorbox}
\section[\texorpdfstring{\hi{इच्छा} (icchā)}{इच्छा (icchā)}]{\hi{इच्छा} (icchā) — a wish, a desire}

\label{lesson:HI-C390-iccha}



\begin{quote}
Before the new one: say the Hindi for to brighten, then the Hindi for a goal.
\end{quote}

\subsection*{You'll want to know: \hi{इच्छा}}
\textbf{\hi{इच्छा}} — \emph{icchā} — \textquotedblleft{}a wish, a desire\textquotedblright{}.

\begin{quote}
\hi{अगर} \hi{तुम्हारी} \hi{इच्छा} \hi{होगी} \hi{तो} \hi{हम} \hi{कल} \hi{पहाड़} \hi{पर} \hi{चलेंगे।}
\end{quote}

\begin{quote}
\emph{agar tumhārī icchā hogī to ham kal pahāṛ par caleṅge.}
\end{quote}

\begin{quote}
If you'd like to, we'll go up the mountain tomorrow.
\end{quote}

Feminine.

\begin{cousinweb}
\textbf{\hi{इच्छा}} (\emph{icchā}) is a tatsama, borrowed from Sanskrit \emph{icchā}, \textquotedblleft{}a wish, a desire\textquotedblright{} (Monier-Williams).
\end{cousinweb}

\begin{grammarlens}[title={What you've built}]
One more word for saying what happens if.
\end{grammarlens}

\subsection*{Your turn}
\begin{itemize}
\raggedright
  \item \emph{Say it:} \emph{icchā}
  \item \emph{Say it:} \emph{icchā}, once more
  \item \emph{Say it:} \emph{agar tumhārī icchā hogī to ham kal pahāṛ par caleṅge.} — If you'd like to, we'll go up the mountain tomorrow.
  \item \emph{Recall:} say the Hindi for to brighten, then the Hindi for a goal, then say \emph{icchā} again
\end{itemize}

\begin{tcolorbox}[breakable,colback=teal!4,colframe=teal!35!black,title={Before you move on}]
What does \hi{इच्छा} mean? (\textquotedblleft{}A wish, a desire\textquotedblright{}.) Say it once more.
\end{tcolorbox}
\section[\texorpdfstring{\hi{संकल्प} (saṅkalp)}{संकल्प (saṅkalp)}]{\hi{संकल्प} (saṅkalp) — a resolve, a firm decision}

\label{lesson:HI-C390-sankalp}



\begin{quote}
Before the new one: say the Hindi for a goal, then the Hindi for a wish.
\end{quote}

\subsection*{You'll want to know: \hi{संकल्प}}
\textbf{\hi{संकल्प}} — \emph{saṅkalp} — \textquotedblleft{}a resolve, a firm decision\textquotedblright{}.

\begin{quote}
\hi{अगर} \hi{संकल्प} \hi{पक्का} \hi{होगा} \hi{तो} \hi{कोई} \hi{तुम्हें} \hi{रोक} \hi{नहीं} \hi{सकेगा।}
\end{quote}

\begin{quote}
\emph{agar saṅkalp pakkā hogā to koī tumhẽ rok nahīṃ sakegā.}
\end{quote}

\begin{quote}
If your resolve is firm, nobody will be able to stop you.
\end{quote}

Masculine.

\begin{cousinweb}
\textbf{\hi{संकल्प}} (\emph{saṅkalp}) is a tatsama, borrowed from Sanskrit \emph{saṃkalpa}, \textquotedblleft{}a notion; a purpose, a will\textquotedblright{} (Monier-Williams).
\end{cousinweb}

\begin{grammarlens}[title={What you've built}]
One more word for saying what happens if.
\end{grammarlens}

\subsection*{Your turn}
\begin{itemize}
\raggedright
  \item \emph{Say it:} \emph{saṅkalp}
  \item \emph{Say it:} \emph{saṅkalp}, once more
  \item \emph{Say it:} \emph{agar saṅkalp pakkā hogā to koī tumhẽ rok nahīṃ sakegā.} — If your resolve is firm, nobody will be able to stop you.
  \item \emph{Recall:} say the Hindi for a goal, then the Hindi for a wish, then say \emph{saṅkalp} again
\end{itemize}

\begin{tcolorbox}[breakable,colback=teal!4,colframe=teal!35!black,title={Before you move on}]
What does \hi{संकल्प} mean? (\textquotedblleft{}A resolve, a firm decision\textquotedblright{}.) Say it once more.
\end{tcolorbox}
\section[\texorpdfstring{\hi{क़दम} (qadam)}{क़दम (qadam)}]{\hi{क़दम} (qadam) — a step}

\label{lesson:HI-C390-qadam}



\begin{quote}
Before the new one: say the Hindi for a wish, then the Hindi for a resolve.
\end{quote}

\subsection*{You'll want to know: \hi{क़दम}}
\textbf{\hi{क़दम}} — \emph{qadam} — \textquotedblleft{}a step\textquotedblright{}.

\begin{quote}
\hi{अगर} \hi{हम} \hi{हर} \hi{दिन} \hi{एक} \hi{क़दम} \hi{आगे} \hi{बढ़ेंगे} \hi{तो} \hi{एक} \hi{दिन} \hi{ज़रूर} \hi{पहुँचेंगे।}
\end{quote}

\begin{quote}
\emph{agar ham har din ek qadam āge baṛheṅge to ek din zarūr pahuñceṅge.}
\end{quote}

\begin{quote}
If we move one step forward every day, we'll get there one day.
\end{quote}

Masculine.

\begin{cousinweb}
\textbf{\hi{क़दम}} (\emph{qadam}) is usually traced to Arabic \emph{qadam}, \textquotedblleft{}a foot, a step\textquotedblright{}, which reached Hindi through Persian.
\end{cousinweb}

\begin{grammarlens}[title={What you've built}]
One more word for saying what happens if.
\end{grammarlens}

\subsection*{Your turn}
\begin{itemize}
\raggedright
  \item \emph{Say it:} \emph{qadam}
  \item \emph{Say it:} \emph{qadam}, once more
  \item \emph{Say it:} \emph{agar ham har din ek qadam āge baṛheṅge to ek din zarūr pahuñceṅge.} — If we move one step forward every day, we'll get there one day.
  \item \emph{Recall:} say the Hindi for a wish, then the Hindi for a resolve, then say \emph{qadam} again
\end{itemize}

\begin{tcolorbox}[breakable,colback=teal!4,colframe=teal!35!black,title={Before you move on}]
What does \hi{क़दम} mean? (\textquotedblleft{}A step\textquotedblright{}.) Say it once more.
\end{tcolorbox}
\section[\texorpdfstring{\hi{कोशिश} \hi{करना} (kośiś karnā)}{कोशिश करना (kośiś karnā)}]{\hi{कोशिश} \hi{करना} (kośiś karnā) — to try, to make an attempt}

\label{lesson:HI-C390-kosiskarna}



\begin{quote}
Before the new one: say the Hindi for a resolve, then the Hindi for a step.
\end{quote}

\subsection*{You'll want to know: \hi{कोशिश} \hi{करना}}
\textbf{\hi{कोशिश} \hi{करना}} — \emph{kośiś karnā} — \textquotedblleft{}to try, to make an attempt\textquotedblright{}.

\begin{quote}
\hi{अगर} \hi{तुम} \hi{कोशिश} \hi{करोगे} \hi{तो} \hi{ज़रूर} \hi{कामयाब} \hi{होगे।}
\end{quote}

\begin{quote}
\emph{agar tum kośiś karoge to zarūr kāmyāb hoge.}
\end{quote}

\begin{quote}
If you try, you'll certainly succeed.
\end{quote}

Past, with \hi{ने} (he or she …): \hi{उसने} \hi{बहुत} \hi{कोशिश} \hi{की}, \emph{usne bahut kośiś kī}.

\begin{cousinweb}
\textbf{\hi{कोशिश} \hi{करना}} (\emph{kośiś karnā}) joins \hi{कोशिश} to \hi{करना}, \textquotedblleft{}to do\textquotedblright{}. \hi{कोशिश} is usually traced to Persian \emph{kōshish}, \textquotedblleft{}an effort\textquotedblright{}.
\end{cousinweb}

\begin{grammarlens}[title={What you've built}]
One more word for saying what happens if.
\end{grammarlens}

\subsection*{Your turn}
\begin{itemize}
\raggedright
  \item \emph{Say it:} \emph{kośiś karnā}
  \item \emph{Say it:} \emph{kośiś karnā}, once more
  \item \emph{Say it:} \emph{agar tum kośiś karoge to zarūr kāmyāb hoge.} — If you try, you'll certainly succeed.
  \item \emph{Recall:} say the Hindi for a resolve, then the Hindi for a step, then say \emph{kośiś karnā} again
\end{itemize}

\begin{tcolorbox}[breakable,colback=teal!4,colframe=teal!35!black,title={Before you move on}]
What does \hi{कोशिश} \hi{करना} mean? (\textquotedblleft{}To try, to make an attempt\textquotedblright{}.) Say it once more.
\end{tcolorbox}
\section[(dialogue)]{Looking back: from Danger and Damage to Setting a Goal}

\label{lesson:HI-R390-review-383-390}



\begin{quote}
Say the Hindi for a crisis and damage.
\end{quote}

\subsection*{Your turn}
Say these aloud:

\begin{itemize}
\raggedright
  \item a crisis, damage, deadly, to unravel, to get wet
  \item to crack, to freeze, to come to the boil, to soften, moderate
  \item to strew, to be wiped out, to tell off, to set, unfinished
  \item to slip away, to be missed, someone to lean on, stubbornness, eligible
  \item extremely, extremely (formal), a bit, plenty, tremendous
  \item more (formal), a quantity, sufficient, minor, close to
  \item better, worse, first-rate, normal, to brighten
  \item a goal, a wish, a resolve, a step, to try
\end{itemize}

\begin{tcolorbox}[breakable,colback=teal!4,colframe=teal!35!black,title={Before you move on}]
A crisis? (\textbf{\hi{संकट}}, \emph{saṅkaṭ}.) Damage? (\textbf{\hi{हानि}}, \emph{hāni}.) Deadly? (\textbf{\hi{घातक}}, \emph{ghātak}.) To unravel? (\textbf{\hi{उधेड़ना}}, \emph{udheṛnā}.) To get wet? (\textbf{\hi{भीगना}}, \emph{bhīgnā}.) To crack? (\textbf{\hi{तड़कना}}, \emph{taṛaknā}.) To freeze? (\textbf{\hi{जमना}}, \emph{jamnā}.) To come to the boil? (\textbf{\hi{उबलना}}, \emph{ubalnā}.) To soften? (\textbf{\hi{गलना}}, \emph{galnā}.) Moderate? (\textbf{\hi{मध्यम}}, \emph{madhyam}.) To strew? (\textbf{\hi{छितराना}}, \emph{chitrānā}.) To be wiped out? (\textbf{\hi{मिटना}}, \emph{miṭnā}.) To tell off? (\textbf{\hi{लताड़ना}}, \emph{latāṛnā}.) To set? (\textbf{\hi{ढलना}}, \emph{ḍhalnā}.) Unfinished? (\textbf{\hi{अधूरा}}, \emph{adhūrā}.) To slip away? (\textbf{\hi{खिसकना}}, \emph{khisaknā}.) To be missed? (\textbf{\hi{खलना}}, \emph{khalnā}.) Someone to lean on? (\textbf{\hi{सहारा}}, \emph{sahārā}.) Stubbornness? (\textbf{\hi{ज़िद}}, \emph{zid}.) Eligible? (\textbf{\hi{योग्य}}, \emph{yogya}.) Extremely? (\textbf{\hi{बेहद}}, \emph{behad}.) Extremely (formal)? (\textbf{\hi{अत्यंत}}, \emph{atyant}.) A bit? (\textbf{\hi{ज़रा}}, \emph{zarā}.) Plenty? (\textbf{\hi{ख़ूब}}, \emph{khūb}.) Tremendous? (\textbf{\hi{ज़बरदस्त}}, \emph{zabardast}.) More (formal)? (\textbf{\hi{अधिक}}, \emph{adhik}.) A quantity? (\textbf{\hi{तादाद}}, \emph{tādād}.) Sufficient? (\textbf{\hi{पर्याप्त}}, \emph{paryāpt}.) Minor? (\textbf{\hi{मामूली}}, \emph{māmūlī}.) Close to? (\textbf{\hi{क़रीब}}, \emph{qarīb}.) Better? (\textbf{\hi{बेहतर}}, \emph{behtar}.) Worse? (\textbf{\hi{बदतर}}, \emph{badtar}.) First-rate? (\textbf{\hi{बढ़िया}}, \emph{baṛhiyā}.) Normal? (\textbf{\hi{सामान्य}}, \emph{sāmānya}.) To brighten? (\textbf{\hi{निखरना}}, \emph{nikharnā}.) A goal? (\textbf{\hi{लक्ष्य}}, \emph{lakṣya}.) A wish? (\textbf{\hi{इच्छा}}, \emph{icchā}.) A resolve? (\textbf{\hi{संकल्प}}, \emph{saṅkalp}.) A step? (\textbf{\hi{क़दम}}, \emph{qadam}.) To try? (\textbf{\hi{कोशिश} \hi{करना}}, \emph{kośiś karnā}.)
\end{tcolorbox}
